\chapter{Crédit à forme réduite}\label{ch:rc:reduced-form-credit}

Le 8 avril 2009, le marché nord-américain des credit default swaps a changé ses contrats
du jour au lendemain. Jusqu'alors, chaque transaction portait son propre spread courant, négocié à
la conclusion, et deux swaps sur le même nom avec la même échéance pouvaient
verser des coupons différents jusqu'à leur terme. À partir de ce jour, les contrats sur les noms investment grade versaient un coupon fixe de 100
points de base par an, et ceux sur les noms high yield 500, et la différence entre le
coupon fixe et le spread du marché a été réglée en espèces dès le départ. Pour convertir un
spread coté en ce montant en espèces, il fallait un modèle sur lequel tout le monde s'accorde, et
le secteur en a publié un en open source. Ce modèle est le \omterm{def:rc:reduced-form-credit:reduced}{modèle à forme réduite} de ce chapitre~:
le défaut est le premier saut d'un processus d'intensité, l'intensité de défaut du
livre~2, chapitre~23. Le livre~2 valorisait un default swap avec une unique intensité de défaut constante~; ici
l'intensité de défaut possède une structure par terme, l'actualisation provient d'une véritable courbe, et les
risques d'une position, par rapport au spread comme au défaut lui-même, sont mesurés.

\section{Le défaut comme premier saut d'une intensité}

\begin{definition}[Modèle à forme réduite]\label{def:rc:reduced-form-credit:reduced}
Un \emph{modèle à forme réduite}\index{modèle à forme réduite} de crédit fait du temps de
défaut $\tau_C$ d'un nom le premier saut d'un processus de comptage d'intensité $\lambda_t$ (un
processus de Poisson à taux dépendant du temps ou aléatoire~; un processus de Cox du
livre~4, chapitre~7, lorsque l'intensité est aléatoire), sans modéliser pourquoi l'entreprise fait défaut~:
$\P(\tau_C\in[t,t+dt)\mid\tau_C>t) = \lambda_t\,dt$ et $Q_C(t) = \E[e^{-\int_0^t\lambda_s\,ds}]$.
\end{definition}

\begin{definition}[Probabilité de défaut, perte en cas de défaut]\label{def:rc:reduced-form-credit:pd}
La \emph{probabilité de défaut}\index{probabilité de défaut} (PD) d'un nom sur un horizon
$T$ vaut $1-Q_C(T)$. La \emph{perte en cas de défaut}\index{perte en cas de défaut}
(LGD) est la part d'une exposition perdue en cas de défaut, $\mathrm{LGD} = 1-R$, où $R$ désigne le
taux de recouvrement.
\end{definition}

\begin{definition}[Facteur d'actualisation risqué, annuité risquée]\label{def:rc:reduced-form-credit:risky}
Le \emph{facteur d'actualisation risqué}\index{facteur d'actualisation risqué} d'un nom à $T$ vaut
$P(0,T)Q_C(T)$, soit la valeur d'une unité versée en $T$ si le nom a survécu, avec
des intensités de défaut déterministes et indépendantes des taux d'intérêt. L'\emph{annuité
risquée}\index{annuité risquée} d'un default swap est la valeur d'une unité par an de prime
versée tant que le nom survit, $\sum_i\delta_iP(0,t_i)Q_C(t_i)$, augmentée de la prime
courue jusqu'à un défaut survenant à l'intérieur de chaque période.
\end{definition}

\begin{proposition}[Le triangle du crédit]\label{prop:rc:reduced-form-credit:triangle}
Avec une intensité de défaut constante $\lambda$, un paiement continu de la prime et un recouvrement $R$ versé au
défaut, le spread au pair d'un default swap, quelle que soit son échéance, vaut
$\mathcal S = \lambda(1-R)$.
\end{proposition}
\begin{proof}
Jambe de prime $\mathcal S\int_0^Te^{-(r+\lambda)t}dt$~; jambe de protection
$(1-R)\int_0^T\lambda e^{-(r+\lambda)t}dt$. Les intégrales sont égales~; le spread au pair
égalise les deux jambes.
\end{proof}

\begin{definition}[Triangle du crédit]\label{def:rc:reduced-form-credit:triangledef}
Le \emph{triangle du crédit}\index{triangle du crédit} est la règle empirique $\mathcal S \approx
\lambda\times\mathrm{LGD}$ de la \cref{prop:rc:reduced-form-credit:triangle}, utilisée pour
lire de tête une intensité de défaut à partir d'un spread (et inversement).
\end{definition}

\section{Valorisation des jambes d'un default swap}

Avec des dates de prime trimestrielles $t_k$ et un défaut dans $(t_{k-1},t_k]$ réglé à $t_k$ avec
une demi-période de prime courue, les deux jambes, sur une \omterm{def:rc:multi-curve-and-collateral-discounting:curves}{courbe d'actualisation} $P$ et une
courbe de survie $Q_C$ valent
\begin{align*}
\text{prime} &= \mathcal S\sum_k\delta_kP(0,t_k)\Bigl(Q_C(t_k)+\tfrac12\bigl(Q_C(t_{k-1})-Q_C(t_k)\bigr)\Bigr),\\
\text{protection} &= (1-R)\sum_kP(0,t_k)\bigl(Q_C(t_{k-1})-Q_C(t_k)\bigr).
\end{align*}

La demi-période courue représente la prime que l'acheteur doit depuis la dernière date de
paiement jusqu'au défaut~; le règlement en fin de période, au lieu d'un règlement au défaut, constitue une petite
erreur d'actualisation que le modèle standard supprime en intégrant sur les dates de défaut.
La \Cref{fig:rc:reduced-form-credit:legs} montre les deux jambes sous forme de flux de trésorerie.

\begin{omfigure}
\begin{tikzpicture}[x=1.05cm,y=1cm,font=\small]
\draw[->,thick] (0,0) -- (10.6,0) node[right] {$t$};
\foreach \k in {1,...,10} {\draw (\k,0.08) -- (\k,-0.08);}
\node[below] at (0,-0.1) {$0$};
\node[below] at (10,-0.1) {$T$};
\foreach \k in {1,...,6} {\draw[->,omDef,thick] (\k,0) -- (\k,-0.9);}
\node[omDef,align=left,anchor=west] at (0.1,-1.35) {prime $\mathcal S\,\delta$ à chaque trimestre tant que le nom survit};
\draw[omThm,thick,dashed] (6.45,0.9) -- (6.45,-0.2);
\node[omThm,anchor=east] at (6.35,0.7) {défaut $\tau_C$};
\draw[->,omThm,very thick] (7,0) -- (7,1.6);
\node[omThm,align=left,anchor=west] at (7.15,1.3) {protection $1-R$ versée,\\ prime courue due};
\draw[omExo,dotted,thick] (7.2,-0.9) -- (10,-0.9);
\node[omExo,anchor=west] at (7.25,-0.6) {aucune prime ultérieure};
\end{tikzpicture}

\omcaption{Les deux jambes d'un default swap. L'acheteur verse la prime trimestriellement (en bleu)
jusqu'au défaut~; au défaut, le vendeur verse la \omterm{def:rc:reduced-form-credit:pd}{perte en cas de défaut} et l'acheteur la prime
courue depuis le dernier paiement (en rouge). Schéma.}\label{fig:rc:reduced-form-credit:legs}
\end{omfigure}

\omcode{code/firm/cdscurve/firm_cdscurve.py}{40}{48}{L'annuité risquée et la jambe de protection sur une courbe d'actualisation et une courbe d'intensité de défaut.}\label{lst:rc:reduced-form-credit:legs}

\section{Bootstrap d'une courbe d'intensité de défaut}

\begin{method}[Courbe d'intensité de défaut à partir des spreads au pair]\label{met:rc:reduced-form-credit:boot}
Prendre les spreads au pair pour des échéances croissantes $T_1<T_2<\dots$ et un taux de recouvrement. Résoudre pour
une intensité de défaut constante sur $(0,T_1]$ qui reproduit le premier~; en la conservant, une intensité de défaut constante sur
$(T_1,T_2]$ qui reproduit le deuxième~; et ainsi de suite. La courbe de survie est
$Q_C(t) = e^{-\int_0^t\lambda}$ avec des intensités de défaut constantes par morceaux. Vérifier que chaque donnée
est repricée et qu'aucune intensité de défaut n'est négative (une intensité négative signifie que les cotations
sont incohérentes).
\end{method}

\omcode{code/firm/cdscurve/firm_cdscurve.py}{56}{68}{Bootstrap d'intensités de défaut constantes par morceaux, de sorte que chaque default swap au pair soit repricé.}\label{lst:rc:reduced-form-credit:boot}

\begin{example}[Une courbe investment grade]\label{ex:rc:reduced-form-credit:curve}
Spreads au pair de 60, 75, 90, 120, 135 et 150 points de base à un, deux, trois, cinq, sept
et dix ans (à titre illustratif), recouvrement de 40~\%, sur la courbe SOFR du chapitre~1, le bootstrap conduisant
à des intensités de défaut de 1.00~\%, 1.51~\%, 2.04~\%, 2.86~\%, 3.03~\% et 3.32~\% sur les intervalles successifs
(\cref{fig:rc:reduced-form-credit:hazard}). Le \omterm{def:rc:reduced-form-credit:triangledef}{triangle du crédit} donne 1.00~\%, 1.25~\%,
1.50~\%, 2.00~\%, 2.25~\% et 2.50~\%~: l'intensité de défaut moyenne jusqu'à chaque échéance, alors que le
bootstrap donne l'intensité de défaut à terme sur chaque intervalle. La \omterm{def:rc:reduced-form-credit:pd}{probabilité de défaut} vaut
0.995~\% sur un an, 9.76~\% sur cinq ans et 23.1~\% sur dix ans
(\cref{fig:rc:reduced-form-credit:survival}).
\end{example}

\begin{omfigure}
\begin{tikzpicture}
\begin{axis}[width=11.5cm,height=5.2cm,
  xlabel={temps (années)},ylabel={intensité de défaut (\%)},
  tick label style={font=\small},label style={font=\small},
  xmin=0,xmax=10,ymin=0,ymax=3.8,grid=major,grid style={gray!25},
  legend columns=2,legend style={font=\footnotesize,at={(0.5,-0.25)},anchor=north,draw=none,fill=none,/tikz/every even column/.append style={column sep=8pt}},
  table/col sep=comma]
\addplot[omThm,very thick,const plot] table[x=t,y=hazard]{figdata/rates-credit-risk/13-reduced-form-credit/hazard.csv};
\addplot[omDef,thick,dashed,const plot] table[x=t,y=triangle]{figdata/rates-credit-risk/13-reduced-form-credit/hazard.csv};
\legend{intensité de défaut bootstrappée (forward),triangle du crédit $\mathcal S/(1-R)$ à chaque échéance}
\end{axis}
\end{tikzpicture}

\omcaption{La courbe d'intensité de défaut bootstrappée à partir de six spreads au pair, et l'intensité moyenne
que lit le \omterm{def:rc:reduced-form-credit:triangledef}{triangle du crédit} à partir de chaque spread. Une courbe de spreads croissante signifie
des intensités forward plus élevées aux échéances lointaines. Données~: spreads illustratifs sur la courbe du chapitre~1~;
le tutoriel du chapitre.}\label{fig:rc:reduced-form-credit:hazard}
\end{omfigure}

\begin{omfigure}
\begin{tikzpicture}
\begin{axis}[width=11.5cm,height=5.2cm,
  xlabel={horizon (années)},ylabel={probabilité (\%)},
  tick label style={font=\small},label style={font=\small},
  xmin=0,xmax=10,ymin=0,ymax=100,grid=major,grid style={gray!25},
  legend columns=2,legend style={font=\footnotesize,at={(0.5,-0.25)},anchor=north,draw=none,fill=none,/tikz/every even column/.append style={column sep=8pt}},
  table/col sep=comma]
\addplot[omDef,very thick] table[x=t,y=q]{figdata/rates-credit-risk/13-reduced-form-credit/survival.csv};
\addplot[omThm,very thick,dashed] table[x=t,y=pd]{figdata/rates-credit-risk/13-reduced-form-credit/survival.csv};
\legend{probabilité de survie $Q_C(t)$,probabilité de défaut $1-Q_C(t)$}
\end{axis}
\end{tikzpicture}

\omcaption{Probabilités de survie et de défaut de l'émetteur bootstrappé jusqu'à dix ans. La
courbe de survie fléchit d'autant plus vite que l'intensité forward augmente~; la \omterm{def:rc:reduced-form-credit:pd}{probabilité de défaut}
atteint 9.8~\% à cinq ans et 23.1~\% à dix ans. Données~: le tutoriel du chapitre.}\label{fig:rc:reduced-form-credit:survival}
\end{omfigure}

\begin{remark}[Le recouvrement est une convention, la perte attendue ne l'est pas]\label{rem:rc:reduced-form-credit:recovery}
Les spreads déterminent la perte attendue, à peu près $\lambda(1-R)$, et non l'intensité seule. Avec un
taux de recouvrement de 25~\% au lieu de 40~\%, les mêmes spreads bootstrappent des intensités de 0.80~\% à 2.64~\%
et une \omterm{def:rc:reduced-form-credit:pd}{probabilité de défaut} à cinq ans de 7.87~\% au lieu de 9.76~\%, alors que tout swap coté
vaut la même chose. Le recouvrement compte pour les positions qui ne sont pas des swaps au pair~: une obligation
cotée proche du pair, un swap éloigné de son coupon, et le \omterm{def:rc:reduced-form-credit:jtd}{risque de saut au défaut} de chacun d'entre eux.
\end{remark}

\begin{remark}[Cotations incohérentes]\label{rem:rc:reduced-form-credit:negative}
Si le spread à trois ans valait 50 points de base, inférieur au spread à deux ans de 75, la revalorisation
de ce spread exigerait une intensité négative la troisième année~: la bissection du \cref{lst:rc:reduced-form-credit:boot}
s'arrête à zéro et ne reproduit que 51.0 points de base. Une courbe qui échoue à ce contrôle
signale au desk qu'une cotation est périmée, ou que le taux de recouvrement retenu est erroné.
\end{remark}

\section{Le modèle standard et ses conventions}

\begin{definition}[Modèle standard de l'ISDA]\label{def:rc:reduced-form-credit:isda}
Le \emph{modèle standard de l'ISDA}\index{modèle standard de l'ISDA} est l'outil de valorisation à forme réduite en source ouverte,
avec des conventions convenues (une intensité de défaut plate implicite d'un spread coté, un taux de recouvrement standard,
la \omterm{def:rc:multi-curve-and-collateral-discounting:curves}{courbe d'actualisation} standard du jour, les règles de calcul des intérêts courus et de paiement),
que le marché utilise pour convertir entre le spread coté d'un contrat et le
paiement initial (upfront) à son coupon standard.
\end{definition}

Le spread coté d'un contrat standard n'est donc pas un spread au pair d'une courbe bootstrappée~:
c'est le spread plat qui, dans le modèle standard, donne le paiement initial négocié. Les deux
coïncident pour l'échéance cotée et diffèrent pour les positions d'autres échéances, que le
desk valorise sur sa courbe bootstrappée.

\begin{example}[Paiement initial à un coupon standard]\label{ex:rc:reduced-form-credit:upfront}
Le contrat à cinq ans coté 120 points de base avec un coupon de 100 points de base coûte
à l'acheteur de protection un paiement initial de 0.872~\% du notionnel (intensité de défaut plate tirée de la cotation, taux zéro
à cinq ans de 3.42~\% plat, taux de recouvrement de 40~\%). Un émetteur coté 500 n'a pas de paiement initial au coupon de 500 et
un paiement initial de 15.06~\% au coupon de 100 (\cref{fig:rc:reduced-form-credit:upfront}).
\end{example}

\begin{omfigure}
\begin{tikzpicture}
\begin{axis}[width=11.5cm,height=5.2cm,
  xlabel={spread coté (pb)},ylabel={paiement initial (\% du notionnel)},
  tick label style={font=\small},label style={font=\small},
  xmin=0,xmax=800,ymin=-25,ymax=25,grid=major,grid style={gray!25},
  legend columns=2,legend style={font=\footnotesize,at={(0.5,-0.25)},anchor=north,draw=none,fill=none,/tikz/every even column/.append style={column sep=8pt}},
  table/col sep=comma]
\addplot[omDef,very thick] table[x=s,y=u100]{figdata/rates-credit-risk/13-reduced-form-credit/upfront.csv};
\addplot[omThm,very thick,dashed] table[x=s,y=u500]{figdata/rates-credit-risk/13-reduced-form-credit/upfront.csv};
\legend{coupon de 100 pb,coupon de 500 pb}
\end{axis}
\end{tikzpicture}

\omcaption{Paiement initial versé par l'acheteur de protection d'un contrat standard à cinq ans, en fonction
du spread coté, aux deux coupons standard nord-américains. Le paiement initial est nul
lorsque le spread coté égale le coupon, et négatif (reçu) en dessous. Données~: la conversion standard
à intensité de défaut plate~; le tutoriel du chapitre.}\label{fig:rc:reduced-form-credit:upfront}
\end{omfigure}

\begin{dated}{2026-09}{Contrats standard}\label{dat:rc:reduced-form-credit:standard}
Depuis le Big Bang Protocol du 8 avril 2009, les credit default swaps (CDS) mono-émetteur nord-américains
se négocient avec un coupon de 100 points de base pour les émetteurs investment grade et de 500 pour les émetteurs
high yield, avec un paiement initial compensant la différence, et les événements de crédit se règlent par enchère.
Le modèle standard utilisé pour la conversion est publié comme code en source ouverte sous une
licence de l'ISDA et administré par S\&P Global Market Intelligence.
\end{dated}

\section{Obligations sur la courbe d'intensité de défaut}

La même courbe de survie valorise une obligation à coupon fixe de l'émetteur~: chaque coupon est un
\omterm{def:rc:reduced-form-credit:risky}{facteur d'actualisation risqué}, et le défaut verse le taux de recouvrement sur la valeur faciale (recouvrement
du pair) à la fin de la période, comme dans \texttt{risky\_bond}. La valeur d'une obligation sur la courbe
d'intensité de défaut bootstrappée à partir des CDS, comparée à son prix de marché, donne
l'écart entre les prix des deux marchés pour un même crédit.

\begin{example}[Une obligation à cinq ans sur la courbe]\label{ex:rc:reduced-form-credit:bond}
Une obligation à cinq ans de 5~\% de l'émetteur (coupons semestriels) vaut 107.08 pour 100 sur la
courbe sans risque et 101.38 sur la courbe d'intensité de défaut bootstrappée. Le \omterm{def:rc:reduced-form-credit:risky}{facteur d'actualisation risqué} à
cinq ans vaut 0.7604, contre 0.8427 sans risque. Si l'obligation se négocie à 97.00, l'intensité
constante qui la réévalue correspond à un spread de CDS de 220 points de base~: l'obligation
est 100 points de base moins chère que ce qu'indique le CDS, soit une base négative de
$-99.8$ points de base (spread de CDS moins spread implicite de l'obligation).
\end{example}

Une base négative incite à acheter l'obligation et à acheter de la protection~: le montage encaisse
le spread excédentaire de l'obligation tandis que la protection couvre le défaut. Ce n'est pas un arbitrage~: il
doit être financé, il est valorisé sur une base susceptible de s'élargir, et son paiement en cas de défaut dépend
du ratio de couverture et du recouvrement~; le problème du week-end le mesure.

\section{Mesures du risque de crédit~: CS01 et saut au défaut}

\begin{definition}[CS01]\label{def:rc:reduced-form-credit:cs01}
La \emph{CS01}\index{CS01} d'une position est la variation de sa valeur quand un spread de
crédit augmente d'un point de base~: par compartiments, par échéance cotée, la courbe
d'intensité étant rebootstrappée, ou parallèle.
\end{definition}

\begin{definition}[Risque de saut au défaut]\label{def:rc:reduced-form-credit:jtd}
Le \emph{risque de saut au défaut}\index{risque de saut au défaut} (jump-to-default, JTD) est la variation de la
valeur d'une position si l'émetteur fait défaut immédiatement~: la protection versée ou reçue,
$\pm(1-R)N$, moins la valeur courante de la position, qui disparaît. Elle n'est pas captée par
\omterm{def:rc:reduced-form-credit:cs01}{CS01}~: le défaut est un saut et non un grand mouvement de spread.
\end{definition}

\begin{example}[Une position de protection]\label{ex:rc:reduced-form-credit:position}
Position longue de protection à cinq ans pour 10~millions d'USD, au coupon de 100 sur la courbe de
l'\cref{ex:rc:reduced-form-credit:curve}~: valant 88\,133~USD (l'\omterm{def:rc:reduced-form-credit:risky}{annuité risquée} vaut 4.41),
avec un \omterm{def:rc:reduced-form-credit:cs01}{CS01} de 4\,391~USD sur la cotation à cinq ans et presque rien ailleurs, et un
gain de saut au défaut de 5.91~millions d'USD.
\end{example}

\begin{definition}[Prime de risque de défaut]\label{def:rc:reduced-form-credit:premium}
La \emph{prime de risque de défaut}\index{prime de risque de défaut} est l'écart entre
la \omterm{def:rc:reduced-form-credit:pd}{probabilité de défaut} risque-neutre implicite dans les spreads et la probabilité réelle
estimée à partir des historiques de défauts~: les spreads rémunèrent les pertes espérées, leur
incertitude et l'illiquidité du crédit, si bien que les PD implicites dépassent les PD historiques,
surtout pour les émetteurs de haute qualité.
\end{definition}

\begin{remark}[Quelle PD utiliser]\label{rem:rc:reduced-form-credit:which}
La valorisation et la couverture utilisent la PD risque-neutre (9.76~\% sur cinq ans, ci-dessus)~; les modèles
de capital, les provisions et les notations utilisent les PD réelles. Si la PD réelle à cinq ans de
l'émetteur était de 1.5~\%, les spreads valoriseraient plus de six fois les pertes espérées liées au défaut. Hull,
Predescu et White (2005) ont trouvé des intensités de défaut risque-neutres à sept ans 16.8 fois
les intensités historiques pour les obligations Aaa, 9.8 pour A, 5.1 pour Baa, 2.1 pour Ba et 1.2 pour B.
\end{remark}

\section{Tutoriel~: une courbe d'intensité de défaut}\label{tut:rc:reduced-form-credit:tutorial}

\begin{tutorial}
\textbf{Objectif.} Bootstrapper une courbe d'intensité à partir des spreads au pair, convertir les cotations en paiements initiaux (upfronts), et
mesurer le \omterm{def:rc:reduced-form-credit:cs01}{CS01} d'une position et son saut au défaut. \textbf{État final~:}
les \cref{fig:rc:reduced-form-credit:hazard,fig:rc:reduced-form-credit:upfront} et les
résultats des \cref{ex:rc:reduced-form-credit:curve,ex:rc:reduced-form-credit:position}.
\begin{tutsteps}
\item \textbf{Bootstrap}~: \texttt{bootstrap(disc, tenors, spreads, recovery)}~; vérifier que chaque
spread au pair est bien réévalué.
\item \textbf{Triangle}~: \texttt{hazard\_table()} contre $\mathcal S/(1-R)$.
\item \textbf{Upfronts}~: \texttt{standard\_upfront} (conversion à intensité constante du livre~2).
\item \textbf{Risque}~: \texttt{protection\_position()}~; \texttt{fig\_rc\_credit.py} écrit les
graphiques.
\end{tutsteps}
\textbf{Pour aller plus loin.} Prendre un recouvrement de 25~\% et comparer les intensités et les upfronts~;
porter le spread à trois ans à 50 points de base et observer le bootstrap échouer à
le réévaluer (la \cref{rem:rc:reduced-form-credit:negative}).
\end{tutorial}

\section{Construction~: la courbe de crédit}\label{bld:rc:reduced-form-credit:cdscurve}

\begin{build}
\bfield{Rôle} Les courbes de crédit de la firme miniature~: chaque CDS, chaque obligation risquée
et chaque contrepartie des chapitres 14 à 20 lit une courbe d'intensité de défaut construite ici.
\bfield{Interface} \texttt{HazardCurve(times, hazards)} avec \texttt{cum} et
\texttt{survival}~; \texttt{legs}, \texttt{par\_spread}, \texttt{bootstrap}~;
\texttt{standard\_upfront}~; \texttt{cds\_value}~; \texttt{risky\_bond}~;
\texttt{bucketed\_cs01}~; \texttt{jump\_to\_default}.
\bfield{Règles} Primes trimestrielles, avec courus d'une demi-période en cas de défaut, défauts réglés en
fin de période~; toute \omterm{def:rc:multi-curve-and-collateral-discounting:curves}{courbe d'actualisation} avec \texttt{df\_t}~; le \texttt{firm\_cds} du livre~2 pour la
conversion à intensité constante.
\bfield{Tests d'acceptation} \texttt{code/firm/cdscurve/tests/}~: chaque cotation est réévaluée~; une
échéance donne l'intensité constante du livre~2~; l'upfront s'annule au coupon~; une obligation
d'intensité nulle est une obligation sans risque~; le \omterm{def:rc:reduced-form-credit:cs01}{CS01} se concentre sur l'échéance propre de la position~; le saut au défaut d'un
acheteur de protection est $(1-R)N$ moins sa valeur.
\bfield{Pour aller plus loin} Calendriers en jours ouvrés et conventions exactes du modèle standard
(step-in, protection à compter de la date de transaction)~; risque sur le taux de recouvrement~; intensité stochastique.
\end{build}

\begin{omsources}
\item Paul, Weiss, «~The Big Bang Protocol and a new structural framework for credit
default swaps~», 24 mars 2009.
\item ISDA CDS Standard Model, \texttt{cdsmodel.com}.
\item D.\ O'Kane, \emph{Modelling Single-name and Multi-name Credit Derivatives}, Wiley,
2008, chapitres 3--7.
\item J.\ Hull, M.\ Predescu et A.\ White, «~Bond prices, default probabilities and risk premiums~»,
\emph{Journal of Credit Risk} 1(2), 2005, 53--60.
\end{omsources}

\section{Exercices}

\begin{exercise}[$\star$]\label{exo:rc:reduced-form-credit:1}
Le spread à cinq ans d'un émetteur est de 180 points de base et le recouvrement de 40~\%. Estimer
l'intensité de défaut et la probabilité de survie à cinq ans à l'aide du \omterm{def:rc:reduced-form-credit:triangledef}{triangle du crédit}.
\end{exercise}

\begin{exercise}[$\star$]\label{exo:rc:reduced-form-credit:2}
Pourquoi l'intensité de défaut bootstrappée est-elle supérieure à celle que donne le \omterm{def:rc:reduced-form-credit:triangledef}{triangle du crédit}
lorsque la courbe des spreads est croissante~?
\end{exercise}

\begin{exercise}[$\star$]\label{exo:rc:reduced-form-credit:3}
Dans l'\cref{ex:rc:reduced-form-credit:position}, quel est le jump-to-default
du vendeur de la même protection~?
\end{exercise}

\begin{exercise}[$\star\star$]\label{exo:rc:reduced-form-credit:4}
Estimer le \omterm{def:rc:reduced-form-credit:cs01}{CS01} de la position de protection à partir de son \omterm{def:rc:reduced-form-credit:risky}{annuité risquée}, et comparer avec
4\,391~USD.
\end{exercise}

\begin{exercise}[$\star\star$]\label{exo:rc:reduced-form-credit:5}
Un émetteur high yield est coté à 350 points de base. Donner l'upfront à cinq ans au coupon de 500
avec les conventions du chapitre, préciser qui le paie, puis donner l'upfront au
coupon de 100.
\end{exercise}

\begin{exercise}[$\star\star$]\label{exo:rc:reduced-form-credit:6}
Pourquoi le jump-to-default n'est-il pas égal au \omterm{def:rc:reduced-form-credit:cs01}{CS01} multiplié par une très grande variation du spread~?
\end{exercise}

\begin{exercise}[$\star\star\star$]\label{exo:rc:reduced-form-credit:7}
\emph{Programmation.} Calculer le package à base négative du problème du week-end~: le
spread implicite de l'obligation, la base, et le jump-to-default du package avec une
couverture par-pour-par.
\end{exercise}

\begin{exercise}[$\star\star\star$]\label{exo:rc:reduced-form-credit:8}
\emph{Trouver la faille.} «~Le spread à cinq ans est de 120 points de base, donc le marché
estime que l'émetteur a une probabilité de 2~\% de faire défaut chaque année.~»
\end{exercise}

\section{Problème~: le package à base négative}

\begin{problem}[{Problème du week-end --- achat d'une obligation bon marché avec protection}]\label{pb:rc:reduced-form-credit:1}
L'émetteur du chapitre possède une obligation à cinq ans à 5~\% (coupon semestriel) qui se traite à 97.00 alors que son
credit default swap (CDS) à cinq ans est coté 120 points de base. Un fonds achète 100~millions d'USD
de nominal de l'obligation et une protection à cinq ans au coupon de 100.

\medskip\noindent\textbf{Partie I --- La base.}
\begin{enumerate}
\item Donner la valeur de l'obligation sur la courbe d'intensité de défaut bootstrappée.
\item Donner le spread au pair plat qui reprice l'obligation à 97.00.
\item Donner la base (spread du credit default swap moins spread implicite de l'obligation).
\item Pourquoi une obligation peut-elle se traiter si loin en dessous de sa valeur implicite dans le credit default swap~?
\item Que gagne le fonds si rien ne se produit~?
\end{enumerate}

\medskip\noindent\textbf{Partie II --- Ratios de couverture.}
\begin{enumerate}[resume]
\item Donner le jump-to-default avec 100~millions d'USD de protection (par-pour-par).
\item Le donner avec 97~millions d'USD (valeur de marché).
\item Donner le \omterm{def:rc:reduced-form-credit:cs01}{CS01} du package dans chaque cas.
\item Donner le taux de recouvrement pour lequel le package en valeur de marché est à l'équilibre en cas de défaut.
\item Quelle couverture retenir, et pourquoi~?
\end{enumerate}

\medskip\noindent\textbf{Partie III --- Risques.}
\begin{enumerate}[resume]
\item Que se passe-t-il pour le package si la base s'élargit encore~?
\item Quel risque de financement porte le fonds~?
\item Quelle valeur représente pour le fonds l'option moins chère à livrer du credit default swap~?
\item Quel risque de contrepartie porte la protection~?
\item Comment le package se comporterait-il lors d'une ruée vers la liquidité (dash for cash) à la manière de 2008~?
\end{enumerate}

\medskip\noindent\textbf{Partie IV --- Jugement.}
\begin{enumerate}[resume]
\item Pourquoi la base négative est-elle appelée un arbitrage, et pourquoi n'en est-ce pas un~?
\item Quels desks l'exploitent, et avec quel bilan~?
\item Que faudrait-il demander avant de la mettre en place~?
\item Énoncer le \emph{résultat nommé}~: la base, ainsi que le jump-to-default et le \omterm{def:rc:reduced-form-credit:cs01}{CS01} du package
avec la couverture par-pour-par.
\item En une phrase~: que mesure une base négative~?
\end{enumerate}
\end{problem}

\section{Questions d'entretien}

\begin{interviewq}[$\star$ \iqroles{trader, chercheur}]\label{iq:rc:reduced-form-credit:1}
Énoncer le \omterm{def:rc:reduced-form-credit:triangledef}{triangle du crédit} et l'appliquer à un spread de 200 points de base.
\end{interviewq}

\begin{interviewq}[$\star\star$ \iqroles{chercheur, développeur}]\label{iq:rc:reduced-form-credit:2}
Comment bootstrappez-vous une courbe d'intensité de défaut, et que peut-il mal se passer~?
\end{interviewq}

\begin{interviewq}[$\star\star$ \iqroles{trader, banque}]\label{iq:rc:reduced-form-credit:3}
Qu'est-ce qui a changé avec la standardisation des credit default swaps en 2009, et pourquoi a-t-elle nécessité un
modèle commun~?
\end{interviewq}

\begin{interviewq}[$\star\star$ \iqroles{risque}]\label{iq:rc:reduced-form-credit:4}
Pourquoi publier le jump-to-default séparément du \omterm{def:rc:reduced-form-credit:cs01}{CS01}~?
\end{interviewq}

\begin{interviewq}[$\star\star\star$ \iqroles{chercheur, risque}]\label{iq:rc:reduced-form-credit:5}
Les \omterm{def:rc:reduced-form-credit:pd}{probabilités de défaut} risque-neutres sont plusieurs fois supérieures aux probabilités historiques. Expliquer, et dire
laquelle vous utiliseriez pour la valorisation, le provisionnement et les fonds propres.
\end{interviewq}

\begin{interviewq}[$\star\star\star$ \iqroles{trader}]\label{iq:rc:reduced-form-credit:6}
Une obligation se traite 100 points de base au-dessus de son credit default swap. Comment la traitez-vous, et que
pourrait-il mal se passer~?
\end{interviewq}
